% Options for packages loaded elsewhere
% Options for packages loaded elsewhere
\PassOptionsToPackage{unicode}{hyperref}
\PassOptionsToPackage{hyphens}{url}
\PassOptionsToPackage{dvipsnames,svgnames,x11names}{xcolor}
%
\documentclass[
  russian,
]{article}
\usepackage{xcolor}
\usepackage[top=20mm,left=20mm,right=20mm,bottom=20mm]{geometry}
\usepackage{amsmath,amssymb}
\setcounter{secnumdepth}{5}
\usepackage{iftex}
\ifPDFTeX
  \usepackage[T1]{fontenc}
  \usepackage[utf8]{inputenc}
  \usepackage{textcomp} % provide euro and other symbols
\else % if luatex or xetex
  \usepackage{unicode-math} % this also loads fontspec
  \defaultfontfeatures{Scale=MatchLowercase}
  \defaultfontfeatures[\rmfamily]{Ligatures=TeX,Scale=1}
\fi
\usepackage{lmodern}
\ifPDFTeX\else
  % xetex/luatex font selection
  \setmainfont[]{Times New Roman}
  \setsansfont[]{Arial}
  \setmonofont[]{Courier New}
\fi
% Use upquote if available, for straight quotes in verbatim environments
\IfFileExists{upquote.sty}{\usepackage{upquote}}{}
\IfFileExists{microtype.sty}{% use microtype if available
  \usepackage[]{microtype}
  \UseMicrotypeSet[protrusion]{basicmath} % disable protrusion for tt fonts
}{}
\makeatletter
\@ifundefined{KOMAClassName}{% if non-KOMA class
  \IfFileExists{parskip.sty}{%
    \usepackage{parskip}
  }{% else
    \setlength{\parindent}{0pt}
    \setlength{\parskip}{6pt plus 2pt minus 1pt}}
}{% if KOMA class
  \KOMAoptions{parskip=half}}
\makeatother
% Make \paragraph and \subparagraph free-standing
\makeatletter
\ifx\paragraph\undefined\else
  \let\oldparagraph\paragraph
  \renewcommand{\paragraph}{
    \@ifstar
      \xxxParagraphStar
      \xxxParagraphNoStar
  }
  \newcommand{\xxxParagraphStar}[1]{\oldparagraph*{#1}\mbox{}}
  \newcommand{\xxxParagraphNoStar}[1]{\oldparagraph{#1}\mbox{}}
\fi
\ifx\subparagraph\undefined\else
  \let\oldsubparagraph\subparagraph
  \renewcommand{\subparagraph}{
    \@ifstar
      \xxxSubParagraphStar
      \xxxSubParagraphNoStar
  }
  \newcommand{\xxxSubParagraphStar}[1]{\oldsubparagraph*{#1}\mbox{}}
  \newcommand{\xxxSubParagraphNoStar}[1]{\oldsubparagraph{#1}\mbox{}}
\fi
\makeatother

\usepackage{color}
\usepackage{fancyvrb}
\newcommand{\VerbBar}{|}
\newcommand{\VERB}{\Verb[commandchars=\\\{\}]}
\DefineVerbatimEnvironment{Highlighting}{Verbatim}{commandchars=\\\{\}}
% Add ',fontsize=\small' for more characters per line
\usepackage{framed}
\definecolor{shadecolor}{RGB}{241,243,245}
\newenvironment{Shaded}{\begin{snugshade}}{\end{snugshade}}
\newcommand{\AlertTok}[1]{\textcolor[rgb]{0.68,0.00,0.00}{#1}}
\newcommand{\AnnotationTok}[1]{\textcolor[rgb]{0.37,0.37,0.37}{#1}}
\newcommand{\AttributeTok}[1]{\textcolor[rgb]{0.40,0.45,0.13}{#1}}
\newcommand{\BaseNTok}[1]{\textcolor[rgb]{0.68,0.00,0.00}{#1}}
\newcommand{\BuiltInTok}[1]{\textcolor[rgb]{0.00,0.23,0.31}{#1}}
\newcommand{\CharTok}[1]{\textcolor[rgb]{0.13,0.47,0.30}{#1}}
\newcommand{\CommentTok}[1]{\textcolor[rgb]{0.37,0.37,0.37}{#1}}
\newcommand{\CommentVarTok}[1]{\textcolor[rgb]{0.37,0.37,0.37}{\textit{#1}}}
\newcommand{\ConstantTok}[1]{\textcolor[rgb]{0.56,0.35,0.01}{#1}}
\newcommand{\ControlFlowTok}[1]{\textcolor[rgb]{0.00,0.23,0.31}{\textbf{#1}}}
\newcommand{\DataTypeTok}[1]{\textcolor[rgb]{0.68,0.00,0.00}{#1}}
\newcommand{\DecValTok}[1]{\textcolor[rgb]{0.68,0.00,0.00}{#1}}
\newcommand{\DocumentationTok}[1]{\textcolor[rgb]{0.37,0.37,0.37}{\textit{#1}}}
\newcommand{\ErrorTok}[1]{\textcolor[rgb]{0.68,0.00,0.00}{#1}}
\newcommand{\ExtensionTok}[1]{\textcolor[rgb]{0.00,0.23,0.31}{#1}}
\newcommand{\FloatTok}[1]{\textcolor[rgb]{0.68,0.00,0.00}{#1}}
\newcommand{\FunctionTok}[1]{\textcolor[rgb]{0.28,0.35,0.67}{#1}}
\newcommand{\ImportTok}[1]{\textcolor[rgb]{0.00,0.46,0.62}{#1}}
\newcommand{\InformationTok}[1]{\textcolor[rgb]{0.37,0.37,0.37}{#1}}
\newcommand{\KeywordTok}[1]{\textcolor[rgb]{0.00,0.23,0.31}{\textbf{#1}}}
\newcommand{\NormalTok}[1]{\textcolor[rgb]{0.00,0.23,0.31}{#1}}
\newcommand{\OperatorTok}[1]{\textcolor[rgb]{0.37,0.37,0.37}{#1}}
\newcommand{\OtherTok}[1]{\textcolor[rgb]{0.00,0.23,0.31}{#1}}
\newcommand{\PreprocessorTok}[1]{\textcolor[rgb]{0.68,0.00,0.00}{#1}}
\newcommand{\RegionMarkerTok}[1]{\textcolor[rgb]{0.00,0.23,0.31}{#1}}
\newcommand{\SpecialCharTok}[1]{\textcolor[rgb]{0.37,0.37,0.37}{#1}}
\newcommand{\SpecialStringTok}[1]{\textcolor[rgb]{0.13,0.47,0.30}{#1}}
\newcommand{\StringTok}[1]{\textcolor[rgb]{0.13,0.47,0.30}{#1}}
\newcommand{\VariableTok}[1]{\textcolor[rgb]{0.07,0.07,0.07}{#1}}
\newcommand{\VerbatimStringTok}[1]{\textcolor[rgb]{0.13,0.47,0.30}{#1}}
\newcommand{\WarningTok}[1]{\textcolor[rgb]{0.37,0.37,0.37}{\textit{#1}}}

\usepackage{longtable,booktabs,array}
\newcounter{none} % for unnumbered tables
\usepackage{calc} % for calculating minipage widths
% Correct order of tables after \paragraph or \subparagraph
\usepackage{etoolbox}
\makeatletter
\patchcmd\longtable{\par}{\if@noskipsec\mbox{}\fi\par}{}{}
\makeatother
% Allow footnotes in longtable head/foot
\IfFileExists{footnotehyper.sty}{\usepackage{footnotehyper}}{\usepackage{footnote}}
\makesavenoteenv{longtable}
\usepackage{graphicx}
\makeatletter
\newsavebox\pandoc@box
\newcommand*\pandocbounded[1]{% scales image to fit in text height/width
  \sbox\pandoc@box{#1}%
  \Gscale@div\@tempa{\textheight}{\dimexpr\ht\pandoc@box+\dp\pandoc@box\relax}%
  \Gscale@div\@tempb{\linewidth}{\wd\pandoc@box}%
  \ifdim\@tempb\p@<\@tempa\p@\let\@tempa\@tempb\fi% select the smaller of both
  \ifdim\@tempa\p@<\p@\scalebox{\@tempa}{\usebox\pandoc@box}%
  \else\usebox{\pandoc@box}%
  \fi%
}
% Set default figure placement to htbp
\def\fps@figure{htbp}
\makeatother



\ifLuaTeX
\usepackage[bidi=basic,shorthands=off,provide=*]{babel}
\else
\usepackage[bidi=default,shorthands=off,provide=*]{babel}
\fi
\ifPDFTeX
\else
\babelfont{rm}[]{Times New Roman}
\fi
\ifLuaTeX
  \usepackage{selnolig} % disable illegal ligatures
\fi


\setlength{\emergencystretch}{3em} % prevent overfull lines

\providecommand{\tightlist}{%
  \setlength{\itemsep}{0pt}\setlength{\parskip}{0pt}}



 


\makeatletter
\@ifpackageloaded{tcolorbox}{}{\usepackage[skins,breakable]{tcolorbox}}
\@ifpackageloaded{fontawesome5}{}{\usepackage{fontawesome5}}
\definecolor{quarto-callout-color}{HTML}{909090}
\definecolor{quarto-callout-note-color}{HTML}{0758E5}
\definecolor{quarto-callout-important-color}{HTML}{CC1914}
\definecolor{quarto-callout-warning-color}{HTML}{EB9113}
\definecolor{quarto-callout-tip-color}{HTML}{00A047}
\definecolor{quarto-callout-caution-color}{HTML}{FC5300}
\definecolor{quarto-callout-color-frame}{HTML}{acacac}
\definecolor{quarto-callout-note-color-frame}{HTML}{4582ec}
\definecolor{quarto-callout-important-color-frame}{HTML}{d9534f}
\definecolor{quarto-callout-warning-color-frame}{HTML}{f0ad4e}
\definecolor{quarto-callout-tip-color-frame}{HTML}{02b875}
\definecolor{quarto-callout-caution-color-frame}{HTML}{fd7e14}
\makeatother
\makeatletter
\@ifpackageloaded{caption}{}{\usepackage{caption}}
\AtBeginDocument{%
\ifdefined\contentsname
  \renewcommand*\contentsname{Содержание}
\else
  \newcommand\contentsname{Содержание}
\fi
\ifdefined\listfigurename
  \renewcommand*\listfigurename{Список Иллюстраций}
\else
  \newcommand\listfigurename{Список Иллюстраций}
\fi
\ifdefined\listtablename
  \renewcommand*\listtablename{Список Таблиц}
\else
  \newcommand\listtablename{Список Таблиц}
\fi
\ifdefined\figurename
  \renewcommand*\figurename{Рисунок}
\else
  \newcommand\figurename{Рисунок}
\fi
\ifdefined\tablename
  \renewcommand*\tablename{Таблица}
\else
  \newcommand\tablename{Таблица}
\fi
}
\@ifpackageloaded{float}{}{\usepackage{float}}
\floatstyle{ruled}
\@ifundefined{c@chapter}{\newfloat{codelisting}{h}{lop}}{\newfloat{codelisting}{h}{lop}[chapter]}
\floatname{codelisting}{Список}
\newcommand*\listoflistings{\listof{codelisting}{Список Каталогов}}
\makeatother
\makeatletter
\makeatother
\makeatletter
\@ifpackageloaded{caption}{}{\usepackage{caption}}
\@ifpackageloaded{subcaption}{}{\usepackage{subcaption}}
\makeatother
\usepackage{bookmark}
\IfFileExists{xurl.sty}{\usepackage{xurl}}{} % add URL line breaks if available
\urlstyle{same}
% fallback for those not using the hyperref driver hyperxmp:
\makeatletter
\@ifundefined{xmpquote}{\newcommand{\xmpquote}[1]{#1}}{}
\makeatother
\hypersetup{
  pdftitle={Задачи по Эконометрике-1: t-тест},
  pdfauthor={Н. В. Артамонов, А. Н. Лата (МГИМО МИД России)},
  pdflang={ru},
  colorlinks=true,
  linkcolor={blue},
  filecolor={Maroon},
  citecolor={Blue},
  urlcolor={Blue},
  pdfcreator={LaTeX via pandoc}}


\title{Задачи по Эконометрике-1: t-тест}
\author{Н. В. Артамонов, А. Н. Лата (МГИМО МИД России)}
\date{}
\begin{document}
\maketitle

\renewcommand*\contentsname{Содержание}
{
\setcounter{tocdepth}{3}
\tableofcontents
}

\begin{tcolorbox}[enhanced jigsaw, arc=.35mm, bottomrule=.15mm, bottomtitle=1mm, breakable, colback=white, colbacktitle=quarto-callout-important-color!10!white, colframe=quarto-callout-important-color-frame, coltitle=black, left=2mm, leftrule=.75mm, opacityback=0, opacitybacktitle=0.6, rightrule=.15mm, title=\textcolor{quarto-callout-important-color}{\faExclamation}\hspace{0.5em}{Важно}, titlerule=0mm, toprule=.15mm, toptitle=1mm]

\begin{itemize}
\tightlist
\item
  Логарифм натуральный; константа включается в модель по умолчанию.
\item
  Данные лежат в папке \texttt{datasets}, описание переменных --- в
  \texttt{datasets/data-description.md}.
\item
  Округление указано в каждой задаче: 3 знака в первом разделе, 2 знака
  во втором. Тестовая статистика считается по неокруглённым оценкам.
\item
  Коэффициенты подписаны так, как их печатает statsmodels:
  \texttt{Intercept} --- константа, \texttt{np.log(x)} --- \(\log x\),
  \texttt{I(np.log(x)\ **\ 2)} --- \(\log^2 x\).
\item
  Индекс коэффициента в гипотезе --- имя переменной: \(\beta_{capital}\)
  --- коэффициент при \texttt{np.log(capital)}. Альтернатива
  двусторонняя.
\item
  Неробастные стандартные ошибки --- обычные ошибки МНК, робастные ---
  HC3; критические значения и P-значения --- по t-распределению
  (\texttt{fit(cov\_type=\textquotesingle{}HC3\textquotesingle{},\ use\_t=True)}).
\end{itemize}

\end{tcolorbox}

\section{t-тест: значимость
коэффициентов}\label{t-ux442ux435ux441ux442-ux437ux43dux430ux447ux438ux43cux43eux441ux442ux44c-ux43aux43eux44dux444ux444ux438ux446ux438ux435ux43dux442ux43eux432}

\subsection{sleep equation \#1}\label{sleep-equation-1}

Для набора данных \texttt{sleep75} рассмотрим регрессию

\texttt{sleep\ \textasciitilde{}\ totwrk\ +\ age\ +\ south\ +\ male\ +\ smsa\ +\ marr\ +\ union}

Переменные \texttt{sleep} и \texttt{totwrk} измерены в минутах в неделю.

\begin{tcolorbox}[enhanced jigsaw, arc=.35mm, bottomrule=.15mm, bottomtitle=1mm, breakable, colback=white, colbacktitle=quarto-callout-important-color!10!white, colframe=quarto-callout-important-color-frame, coltitle=black, left=2mm, leftrule=.75mm, opacityback=0, opacitybacktitle=0.6, rightrule=.15mm, title=\textcolor{quarto-callout-important-color}{\faExclamation}\hspace{0.5em}{Условия теста}, titlerule=0mm, toprule=.15mm, toptitle=1mm]

\begin{itemize}
\tightlist
\item
  уровень значимости: \textbf{5\%}
\item
  стандартные ошибки: \textbf{неробастные}
\item
  гипотеза: \(H_0:\beta_j=0\) для каждого коэффициента
\end{itemize}

\end{tcolorbox}

\begin{itemize}
\tightlist
\item
  Оцените регрессию.
\item
  Выведите результаты t-теста.
\item
  Вычислите критическое значение.
\item
  Какие коэффициенты значимы?
\end{itemize}

\textbf{Ответ округлите до 3-х десятичных знаков.}

\textbf{Результаты \(t\)-теста для коэффициентов (неробастные
стандартные ошибки):}

{\def\LTcaptype{none} % do not increment counter
\begin{longtable}[]{@{}
  >{\raggedright\arraybackslash}p{(\linewidth - 10\tabcolsep) * \real{0.1667}}
  >{\raggedleft\arraybackslash}p{(\linewidth - 10\tabcolsep) * \real{0.1515}}
  >{\raggedleft\arraybackslash}p{(\linewidth - 10\tabcolsep) * \real{0.1818}}
  >{\raggedleft\arraybackslash}p{(\linewidth - 10\tabcolsep) * \real{0.1667}}
  >{\raggedleft\arraybackslash}p{(\linewidth - 10\tabcolsep) * \real{0.1667}}
  >{\raggedright\arraybackslash}p{(\linewidth - 10\tabcolsep) * \real{0.1667}}@{}}
\toprule\noalign{}
\begin{minipage}[b]{\linewidth}\raggedright
\end{minipage} & \begin{minipage}[b]{\linewidth}\raggedleft
Coef.
\end{minipage} & \begin{minipage}[b]{\linewidth}\raggedleft
Std.Err.
\end{minipage} & \begin{minipage}[b]{\linewidth}\raggedleft
t-value
\end{minipage} & \begin{minipage}[b]{\linewidth}\raggedleft
p-value
\end{minipage} & \begin{minipage}[b]{\linewidth}\raggedright
Signif.
\end{minipage} \\
\midrule\noalign{}
\endhead
\bottomrule\noalign{}
\endlastfoot
Intercept & 3441.297 & 78.942 & 43.593 & \textless0.001 & \\
totwrk & -0.169 & 0.018 & -9.363 & \textless0.001 & *** \\
age & 2.841 & 1.387 & 2.048 & 0.041 & ** \\
south & 101.676 & 41.834 & 2.430 & 0.015 & ** \\
male & 86.038 & 34.906 & 2.465 & 0.014 & ** \\
smsa & -53.925 & 33.156 & -1.626 & 0.104 & \\
marr & 30.006 & 41.910 & 0.716 & 0.474 & \\
union & 12.440 & 38.096 & 0.327 & 0.744 & \\
\end{longtable}
}

\emph{Signif. codes:} \texttt{***} \(p<0.01\), \texttt{**} \(p<0.05\),
\texttt{*} \(p<0.1\)

Модель была подогнана по 706 наблюдениям.

\begin{tcolorbox}[enhanced jigsaw, arc=.35mm, bottomrule=.15mm, bottomtitle=1mm, breakable, colback=white, colbacktitle=quarto-callout-note-color!10!white, colframe=quarto-callout-note-color-frame, coltitle=black, left=2mm, leftrule=.75mm, opacityback=0, opacitybacktitle=0.6, rightrule=.15mm, title=\textcolor{quarto-callout-note-color}{\faInfo}\hspace{0.5em}{Критическое значение}, titlerule=0mm, toprule=.15mm, toptitle=1mm]

\(t_{crit} = t_{0.975}(698) = 1.963\)

\end{tcolorbox}

\begin{tcolorbox}[enhanced jigsaw, arc=.35mm, bottomrule=.15mm, bottomtitle=1mm, breakable, colback=white, colbacktitle=quarto-callout-note-color!10!white, colframe=quarto-callout-note-color-frame, coltitle=black, left=2mm, leftrule=.75mm, opacityback=0, opacitybacktitle=0.6, rightrule=.15mm, title=\textcolor{quarto-callout-note-color}{\faInfo}\hspace{0.5em}{Коэффициенты, значимые на уровне 5\% (по неробастным стандартным
ошибкам)}, titlerule=0mm, toprule=.15mm, toptitle=1mm]

\texttt{totwrk}, \texttt{age}, \texttt{south}, \texttt{male}

\end{tcolorbox}

\subsection{sleep equation \#2}\label{sleep-equation-2}

Для набора данных \texttt{sleep75} рассмотрим регрессию

\texttt{sleep\ \textasciitilde{}\ totwrk\ +\ age\ +\ south\ +\ male\ +\ smsa\ +\ marr\ +\ union}

\begin{tcolorbox}[enhanced jigsaw, arc=.35mm, bottomrule=.15mm, bottomtitle=1mm, breakable, colback=white, colbacktitle=quarto-callout-important-color!10!white, colframe=quarto-callout-important-color-frame, coltitle=black, left=2mm, leftrule=.75mm, opacityback=0, opacitybacktitle=0.6, rightrule=.15mm, title=\textcolor{quarto-callout-important-color}{\faExclamation}\hspace{0.5em}{Условия теста}, titlerule=0mm, toprule=.15mm, toptitle=1mm]

\begin{itemize}
\tightlist
\item
  уровень значимости: \textbf{5\%}
\item
  стандартные ошибки: \textbf{робастные (HC3)}
\item
  гипотеза: \(H_0:\beta_j=0\) для каждого коэффициента
\end{itemize}

\end{tcolorbox}

\begin{itemize}
\tightlist
\item
  Оцените регрессию.
\item
  Выведите результаты t-теста.
\item
  Вычислите критическое значение.
\item
  Какие коэффициенты значимы?
\end{itemize}

\textbf{Ответ округлите до 3-х десятичных знаков.}

\textbf{Результаты \(t\)-теста для коэффициентов (робастные стандартные
ошибки, HC3):}

{\def\LTcaptype{none} % do not increment counter
\begin{longtable}[]{@{}
  >{\raggedright\arraybackslash}p{(\linewidth - 10\tabcolsep) * \real{0.1667}}
  >{\raggedleft\arraybackslash}p{(\linewidth - 10\tabcolsep) * \real{0.1515}}
  >{\raggedleft\arraybackslash}p{(\linewidth - 10\tabcolsep) * \real{0.1818}}
  >{\raggedleft\arraybackslash}p{(\linewidth - 10\tabcolsep) * \real{0.1667}}
  >{\raggedleft\arraybackslash}p{(\linewidth - 10\tabcolsep) * \real{0.1667}}
  >{\raggedright\arraybackslash}p{(\linewidth - 10\tabcolsep) * \real{0.1667}}@{}}
\toprule\noalign{}
\begin{minipage}[b]{\linewidth}\raggedright
\end{minipage} & \begin{minipage}[b]{\linewidth}\raggedleft
Coef.
\end{minipage} & \begin{minipage}[b]{\linewidth}\raggedleft
Std.Err.
\end{minipage} & \begin{minipage}[b]{\linewidth}\raggedleft
t-value
\end{minipage} & \begin{minipage}[b]{\linewidth}\raggedleft
p-value
\end{minipage} & \begin{minipage}[b]{\linewidth}\raggedright
Signif.
\end{minipage} \\
\midrule\noalign{}
\endhead
\bottomrule\noalign{}
\endlastfoot
Intercept & 3441.297 & 88.891 & 38.714 & \textless0.001 & \\
totwrk & -0.169 & 0.020 & -8.299 & \textless0.001 & *** \\
age & 2.841 & 1.362 & 2.086 & 0.037 & ** \\
south & 101.676 & 42.315 & 2.403 & 0.017 & ** \\
male & 86.038 & 35.758 & 2.406 & 0.016 & ** \\
smsa & -53.925 & 33.372 & -1.616 & 0.107 & \\
marr & 30.006 & 46.090 & 0.651 & 0.515 & \\
union & 12.440 & 37.277 & 0.334 & 0.739 & \\
\end{longtable}
}

\emph{Signif. codes:} \texttt{***} \(p<0.01\), \texttt{**} \(p<0.05\),
\texttt{*} \(p<0.1\)

Модель была подогнана по 706 наблюдениям.

\begin{tcolorbox}[enhanced jigsaw, arc=.35mm, bottomrule=.15mm, bottomtitle=1mm, breakable, colback=white, colbacktitle=quarto-callout-note-color!10!white, colframe=quarto-callout-note-color-frame, coltitle=black, left=2mm, leftrule=.75mm, opacityback=0, opacitybacktitle=0.6, rightrule=.15mm, title=\textcolor{quarto-callout-note-color}{\faInfo}\hspace{0.5em}{Критическое значение}, titlerule=0mm, toprule=.15mm, toptitle=1mm]

\(t_{crit} = t_{0.975}(698) = 1.963\)

\end{tcolorbox}

\begin{tcolorbox}[enhanced jigsaw, arc=.35mm, bottomrule=.15mm, bottomtitle=1mm, breakable, colback=white, colbacktitle=quarto-callout-note-color!10!white, colframe=quarto-callout-note-color-frame, coltitle=black, left=2mm, leftrule=.75mm, opacityback=0, opacitybacktitle=0.6, rightrule=.15mm, title=\textcolor{quarto-callout-note-color}{\faInfo}\hspace{0.5em}{Коэффициенты, значимые на уровне 5\% (по робастным стандартным ошибкам
HC3)}, titlerule=0mm, toprule=.15mm, toptitle=1mm]

\texttt{totwrk}, \texttt{age}, \texttt{south}, \texttt{male}

\end{tcolorbox}

\subsection{sleep equation \#3}\label{sleep-equation-3}

Для набора данных \texttt{sleep75} рассмотрим регрессию

\texttt{sleep\ \textasciitilde{}\ totwrk\ +\ age\ +\ south\ +\ male\ +\ smsa\ +\ marr\ +\ union\ +\ np.log(hrwage)}

Переменная \texttt{hrwage} известна не для всех наблюдений: модель
подгоняется по 532 наблюдениям из 706.

\begin{tcolorbox}[enhanced jigsaw, arc=.35mm, bottomrule=.15mm, bottomtitle=1mm, breakable, colback=white, colbacktitle=quarto-callout-important-color!10!white, colframe=quarto-callout-important-color-frame, coltitle=black, left=2mm, leftrule=.75mm, opacityback=0, opacitybacktitle=0.6, rightrule=.15mm, title=\textcolor{quarto-callout-important-color}{\faExclamation}\hspace{0.5em}{Условия теста}, titlerule=0mm, toprule=.15mm, toptitle=1mm]

\begin{itemize}
\tightlist
\item
  уровень значимости: \textbf{10\%}
\item
  стандартные ошибки: \textbf{неробастные}
\item
  гипотеза: \(H_0:\beta_j=0\) для каждого коэффициента
\end{itemize}

\end{tcolorbox}

\begin{itemize}
\tightlist
\item
  Оцените регрессию.
\item
  Выведите результаты t-теста.
\item
  Вычислите критическое значение.
\item
  Какие коэффициенты значимы?
\end{itemize}

\textbf{Ответ округлите до 3-х десятичных знаков.}

\textbf{Результаты \(t\)-теста для коэффициентов (неробастные
стандартные ошибки):}

{\def\LTcaptype{none} % do not increment counter
\begin{longtable}[]{@{}
  >{\raggedright\arraybackslash}p{(\linewidth - 10\tabcolsep) * \real{0.2254}}
  >{\raggedleft\arraybackslash}p{(\linewidth - 10\tabcolsep) * \real{0.1408}}
  >{\raggedleft\arraybackslash}p{(\linewidth - 10\tabcolsep) * \real{0.1690}}
  >{\raggedleft\arraybackslash}p{(\linewidth - 10\tabcolsep) * \real{0.1549}}
  >{\raggedleft\arraybackslash}p{(\linewidth - 10\tabcolsep) * \real{0.1549}}
  >{\raggedright\arraybackslash}p{(\linewidth - 10\tabcolsep) * \real{0.1549}}@{}}
\toprule\noalign{}
\begin{minipage}[b]{\linewidth}\raggedright
\end{minipage} & \begin{minipage}[b]{\linewidth}\raggedleft
Coef.
\end{minipage} & \begin{minipage}[b]{\linewidth}\raggedleft
Std.Err.
\end{minipage} & \begin{minipage}[b]{\linewidth}\raggedleft
t-value
\end{minipage} & \begin{minipage}[b]{\linewidth}\raggedleft
p-value
\end{minipage} & \begin{minipage}[b]{\linewidth}\raggedright
Signif.
\end{minipage} \\
\midrule\noalign{}
\endhead
\bottomrule\noalign{}
\endlastfoot
Intercept & 3453.355 & 93.371 & 36.985 & \textless0.001 & \\
totwrk & -0.159 & 0.021 & -7.629 & \textless0.001 & *** \\
age & 1.950 & 1.584 & 1.230 & 0.219 & \\
south & 80.244 & 45.216 & 1.775 & 0.077 & * \\
male & 41.098 & 43.049 & 0.955 & 0.340 & \\
smsa & -36.093 & 38.206 & -0.945 & 0.345 & \\
marr & 58.917 & 46.613 & 1.264 & 0.207 & \\
union & 24.479 & 42.948 & 0.570 & 0.569 & \\
np.log(hrwage) & -2.590 & 32.197 & -0.080 & 0.936 & \\
\end{longtable}
}

\emph{Signif. codes:} \texttt{***} \(p<0.01\), \texttt{**} \(p<0.05\),
\texttt{*} \(p<0.1\)

Модель была подогнана по 532 наблюдениям.

\begin{tcolorbox}[enhanced jigsaw, arc=.35mm, bottomrule=.15mm, bottomtitle=1mm, breakable, colback=white, colbacktitle=quarto-callout-note-color!10!white, colframe=quarto-callout-note-color-frame, coltitle=black, left=2mm, leftrule=.75mm, opacityback=0, opacitybacktitle=0.6, rightrule=.15mm, title=\textcolor{quarto-callout-note-color}{\faInfo}\hspace{0.5em}{Критическое значение}, titlerule=0mm, toprule=.15mm, toptitle=1mm]

\(t_{crit} = t_{0.95}(523) = 1.648\)

\end{tcolorbox}

\begin{tcolorbox}[enhanced jigsaw, arc=.35mm, bottomrule=.15mm, bottomtitle=1mm, breakable, colback=white, colbacktitle=quarto-callout-note-color!10!white, colframe=quarto-callout-note-color-frame, coltitle=black, left=2mm, leftrule=.75mm, opacityback=0, opacitybacktitle=0.6, rightrule=.15mm, title=\textcolor{quarto-callout-note-color}{\faInfo}\hspace{0.5em}{Коэффициенты, значимые на уровне 10\% (по неробастным стандартным
ошибкам)}, titlerule=0mm, toprule=.15mm, toptitle=1mm]

\texttt{totwrk}, \texttt{south}

\end{tcolorbox}

\subsection{sleep equation \#4}\label{sleep-equation-4}

Для набора данных \texttt{sleep75} рассмотрим регрессию

\texttt{sleep\ \textasciitilde{}\ totwrk\ +\ age\ +\ south\ +\ male\ +\ smsa\ +\ marr\ +\ union\ +\ np.log(hrwage)}

Переменная \texttt{hrwage} известна не для всех наблюдений: модель
подгоняется по 532 наблюдениям из 706.

\begin{tcolorbox}[enhanced jigsaw, arc=.35mm, bottomrule=.15mm, bottomtitle=1mm, breakable, colback=white, colbacktitle=quarto-callout-important-color!10!white, colframe=quarto-callout-important-color-frame, coltitle=black, left=2mm, leftrule=.75mm, opacityback=0, opacitybacktitle=0.6, rightrule=.15mm, title=\textcolor{quarto-callout-important-color}{\faExclamation}\hspace{0.5em}{Условия теста}, titlerule=0mm, toprule=.15mm, toptitle=1mm]

\begin{itemize}
\tightlist
\item
  уровень значимости: \textbf{10\%}
\item
  стандартные ошибки: \textbf{робастные (HC3)}
\item
  гипотеза: \(H_0:\beta_j=0\) для каждого коэффициента
\end{itemize}

\end{tcolorbox}

\begin{itemize}
\tightlist
\item
  Оцените регрессию.
\item
  Выведите результаты t-теста.
\item
  Вычислите критическое значение.
\item
  Какие коэффициенты значимы?
\end{itemize}

\textbf{Ответ округлите до 3-х десятичных знаков.}

\textbf{Результаты \(t\)-теста для коэффициентов (робастные стандартные
ошибки, HC3):}

{\def\LTcaptype{none} % do not increment counter
\begin{longtable}[]{@{}
  >{\raggedright\arraybackslash}p{(\linewidth - 10\tabcolsep) * \real{0.2254}}
  >{\raggedleft\arraybackslash}p{(\linewidth - 10\tabcolsep) * \real{0.1408}}
  >{\raggedleft\arraybackslash}p{(\linewidth - 10\tabcolsep) * \real{0.1690}}
  >{\raggedleft\arraybackslash}p{(\linewidth - 10\tabcolsep) * \real{0.1549}}
  >{\raggedleft\arraybackslash}p{(\linewidth - 10\tabcolsep) * \real{0.1549}}
  >{\raggedright\arraybackslash}p{(\linewidth - 10\tabcolsep) * \real{0.1549}}@{}}
\toprule\noalign{}
\begin{minipage}[b]{\linewidth}\raggedright
\end{minipage} & \begin{minipage}[b]{\linewidth}\raggedleft
Coef.
\end{minipage} & \begin{minipage}[b]{\linewidth}\raggedleft
Std.Err.
\end{minipage} & \begin{minipage}[b]{\linewidth}\raggedleft
t-value
\end{minipage} & \begin{minipage}[b]{\linewidth}\raggedleft
p-value
\end{minipage} & \begin{minipage}[b]{\linewidth}\raggedright
Signif.
\end{minipage} \\
\midrule\noalign{}
\endhead
\bottomrule\noalign{}
\endlastfoot
Intercept & 3453.355 & 98.213 & 35.162 & \textless0.001 & \\
totwrk & -0.159 & 0.025 & -6.354 & \textless0.001 & *** \\
age & 1.950 & 1.571 & 1.241 & 0.215 & \\
south & 80.244 & 47.401 & 1.693 & 0.091 & * \\
male & 41.098 & 42.778 & 0.961 & 0.337 & \\
smsa & -36.093 & 38.580 & -0.936 & 0.350 & \\
marr & 58.917 & 48.350 & 1.219 & 0.224 & \\
union & 24.479 & 43.006 & 0.569 & 0.569 & \\
np.log(hrwage) & -2.590 & 32.992 & -0.079 & 0.937 & \\
\end{longtable}
}

\emph{Signif. codes:} \texttt{***} \(p<0.01\), \texttt{**} \(p<0.05\),
\texttt{*} \(p<0.1\)

Модель была подогнана по 532 наблюдениям.

\begin{tcolorbox}[enhanced jigsaw, arc=.35mm, bottomrule=.15mm, bottomtitle=1mm, breakable, colback=white, colbacktitle=quarto-callout-note-color!10!white, colframe=quarto-callout-note-color-frame, coltitle=black, left=2mm, leftrule=.75mm, opacityback=0, opacitybacktitle=0.6, rightrule=.15mm, title=\textcolor{quarto-callout-note-color}{\faInfo}\hspace{0.5em}{Критическое значение}, titlerule=0mm, toprule=.15mm, toptitle=1mm]

\(t_{crit} = t_{0.95}(523) = 1.648\)

\end{tcolorbox}

\begin{tcolorbox}[enhanced jigsaw, arc=.35mm, bottomrule=.15mm, bottomtitle=1mm, breakable, colback=white, colbacktitle=quarto-callout-note-color!10!white, colframe=quarto-callout-note-color-frame, coltitle=black, left=2mm, leftrule=.75mm, opacityback=0, opacitybacktitle=0.6, rightrule=.15mm, title=\textcolor{quarto-callout-note-color}{\faInfo}\hspace{0.5em}{Коэффициенты, значимые на уровне 10\% (по робастным стандартным ошибкам
HC3)}, titlerule=0mm, toprule=.15mm, toptitle=1mm]

\texttt{totwrk}, \texttt{south}

\end{tcolorbox}

\subsection{output equation \#1}\label{output-equation-1}

Для набора данных \texttt{Labour} рассмотрим регрессию

\begin{Shaded}
\begin{Highlighting}[]
\NormalTok{np.log(output) }\OperatorTok{\textasciitilde{}}\NormalTok{ np.log(capital) }\OperatorTok{+}\NormalTok{ np.log(labour) }\OperatorTok{+}\NormalTok{ np.log(wage) }\OperatorTok{+}
\NormalTok{    I(np.log(capital)}\OperatorTok{**}\DecValTok{2}\NormalTok{) }\OperatorTok{+}\NormalTok{ I(np.log(labour)}\OperatorTok{**}\DecValTok{2}\NormalTok{)}
\end{Highlighting}
\end{Shaded}

Переменные \texttt{output} и \texttt{capital} измерены в млн евро,
\texttt{labour} --- число работников, \texttt{wage} --- в тыс. евро на
работника.

\begin{tcolorbox}[enhanced jigsaw, arc=.35mm, bottomrule=.15mm, bottomtitle=1mm, breakable, colback=white, colbacktitle=quarto-callout-important-color!10!white, colframe=quarto-callout-important-color-frame, coltitle=black, left=2mm, leftrule=.75mm, opacityback=0, opacitybacktitle=0.6, rightrule=.15mm, title=\textcolor{quarto-callout-important-color}{\faExclamation}\hspace{0.5em}{Условия теста}, titlerule=0mm, toprule=.15mm, toptitle=1mm]

\begin{itemize}
\tightlist
\item
  уровень значимости: \textbf{10\%}
\item
  стандартные ошибки: \textbf{неробастные}
\item
  гипотеза: \(H_0:\beta_j=0\) для каждого коэффициента
\end{itemize}

\end{tcolorbox}

\begin{itemize}
\tightlist
\item
  Оцените регрессию.
\item
  Выведите результаты t-теста.
\item
  Вычислите критическое значение.
\item
  Какие коэффициенты значимы?
\end{itemize}

\textbf{Ответ округлите до 3-х десятичных знаков.}

\textbf{Результаты \(t\)-теста для коэффициентов (неробастные
стандартные ошибки):}

{\def\LTcaptype{none} % do not increment counter
\begin{longtable}[]{@{}
  >{\raggedright\arraybackslash}p{(\linewidth - 10\tabcolsep) * \real{0.3165}}
  >{\raggedleft\arraybackslash}p{(\linewidth - 10\tabcolsep) * \real{0.1139}}
  >{\raggedleft\arraybackslash}p{(\linewidth - 10\tabcolsep) * \real{0.1519}}
  >{\raggedleft\arraybackslash}p{(\linewidth - 10\tabcolsep) * \real{0.1392}}
  >{\raggedleft\arraybackslash}p{(\linewidth - 10\tabcolsep) * \real{0.1392}}
  >{\raggedright\arraybackslash}p{(\linewidth - 10\tabcolsep) * \real{0.1392}}@{}}
\toprule\noalign{}
\begin{minipage}[b]{\linewidth}\raggedright
\end{minipage} & \begin{minipage}[b]{\linewidth}\raggedleft
Coef.
\end{minipage} & \begin{minipage}[b]{\linewidth}\raggedleft
Std.Err.
\end{minipage} & \begin{minipage}[b]{\linewidth}\raggedleft
t-value
\end{minipage} & \begin{minipage}[b]{\linewidth}\raggedleft
p-value
\end{minipage} & \begin{minipage}[b]{\linewidth}\raggedright
Signif.
\end{minipage} \\
\midrule\noalign{}
\endhead
\bottomrule\noalign{}
\endlastfoot
Intercept & -4.258 & 0.257 & -16.599 & \textless0.001 & \\
np.log(capital) & 0.140 & 0.014 & 9.743 & \textless0.001 & *** \\
np.log(labour) & 0.473 & 0.071 & 6.672 & \textless0.001 & *** \\
np.log(wage) & 0.853 & 0.058 & 14.751 & \textless0.001 & *** \\
I(np.log(capital) ** 2) & 0.008 & 0.004 & 1.735 & 0.083 & * \\
I(np.log(labour) ** 2) & 0.027 & 0.008 & 3.355 & \textless0.001 & *** \\
\end{longtable}
}

\emph{Signif. codes:} \texttt{***} \(p<0.01\), \texttt{**} \(p<0.05\),
\texttt{*} \(p<0.1\)

Модель была подогнана по 569 наблюдениям.

\begin{tcolorbox}[enhanced jigsaw, arc=.35mm, bottomrule=.15mm, bottomtitle=1mm, breakable, colback=white, colbacktitle=quarto-callout-note-color!10!white, colframe=quarto-callout-note-color-frame, coltitle=black, left=2mm, leftrule=.75mm, opacityback=0, opacitybacktitle=0.6, rightrule=.15mm, title=\textcolor{quarto-callout-note-color}{\faInfo}\hspace{0.5em}{Критическое значение}, titlerule=0mm, toprule=.15mm, toptitle=1mm]

\(t_{crit} = t_{0.95}(563) = 1.648\)

\end{tcolorbox}

\begin{tcolorbox}[enhanced jigsaw, arc=.35mm, bottomrule=.15mm, bottomtitle=1mm, breakable, colback=white, colbacktitle=quarto-callout-note-color!10!white, colframe=quarto-callout-note-color-frame, coltitle=black, left=2mm, leftrule=.75mm, opacityback=0, opacitybacktitle=0.6, rightrule=.15mm, title=\textcolor{quarto-callout-note-color}{\faInfo}\hspace{0.5em}{Коэффициенты, значимые на уровне 10\% (по неробастным стандартным
ошибкам)}, titlerule=0mm, toprule=.15mm, toptitle=1mm]

\texttt{np.log(capital)}, \texttt{np.log(labour)},
\texttt{np.log(wage)}, \texttt{I(np.log(capital)\ **\ 2)},
\texttt{I(np.log(labour)\ **\ 2)}

\end{tcolorbox}

\subsection{output equation \#2}\label{output-equation-2}

Для набора данных \texttt{Labour} рассмотрим регрессию

\begin{Shaded}
\begin{Highlighting}[]
\NormalTok{np.log(output) }\OperatorTok{\textasciitilde{}}\NormalTok{ np.log(capital) }\OperatorTok{+}\NormalTok{ np.log(labour) }\OperatorTok{+}\NormalTok{ np.log(wage) }\OperatorTok{+}
\NormalTok{    I(np.log(capital)}\OperatorTok{**}\DecValTok{2}\NormalTok{) }\OperatorTok{+}\NormalTok{ I(np.log(labour)}\OperatorTok{**}\DecValTok{2}\NormalTok{)}
\end{Highlighting}
\end{Shaded}

\begin{tcolorbox}[enhanced jigsaw, arc=.35mm, bottomrule=.15mm, bottomtitle=1mm, breakable, colback=white, colbacktitle=quarto-callout-important-color!10!white, colframe=quarto-callout-important-color-frame, coltitle=black, left=2mm, leftrule=.75mm, opacityback=0, opacitybacktitle=0.6, rightrule=.15mm, title=\textcolor{quarto-callout-important-color}{\faExclamation}\hspace{0.5em}{Условия теста}, titlerule=0mm, toprule=.15mm, toptitle=1mm]

\begin{itemize}
\tightlist
\item
  уровень значимости: \textbf{10\%}
\item
  стандартные ошибки: \textbf{робастные (HC3)}
\item
  гипотеза: \(H_0:\beta_j=0\) для каждого коэффициента
\end{itemize}

\end{tcolorbox}

\begin{itemize}
\tightlist
\item
  Оцените регрессию.
\item
  Выведите результаты t-теста.
\item
  Вычислите критическое значение.
\item
  Какие коэффициенты значимы?
\end{itemize}

\textbf{Ответ округлите до 3-х десятичных знаков.}

\textbf{Результаты \(t\)-теста для коэффициентов (робастные стандартные
ошибки, HC3):}

{\def\LTcaptype{none} % do not increment counter
\begin{longtable}[]{@{}
  >{\raggedright\arraybackslash}p{(\linewidth - 10\tabcolsep) * \real{0.3165}}
  >{\raggedleft\arraybackslash}p{(\linewidth - 10\tabcolsep) * \real{0.1139}}
  >{\raggedleft\arraybackslash}p{(\linewidth - 10\tabcolsep) * \real{0.1519}}
  >{\raggedleft\arraybackslash}p{(\linewidth - 10\tabcolsep) * \real{0.1392}}
  >{\raggedleft\arraybackslash}p{(\linewidth - 10\tabcolsep) * \real{0.1392}}
  >{\raggedright\arraybackslash}p{(\linewidth - 10\tabcolsep) * \real{0.1392}}@{}}
\toprule\noalign{}
\begin{minipage}[b]{\linewidth}\raggedright
\end{minipage} & \begin{minipage}[b]{\linewidth}\raggedleft
Coef.
\end{minipage} & \begin{minipage}[b]{\linewidth}\raggedleft
Std.Err.
\end{minipage} & \begin{minipage}[b]{\linewidth}\raggedleft
t-value
\end{minipage} & \begin{minipage}[b]{\linewidth}\raggedleft
p-value
\end{minipage} & \begin{minipage}[b]{\linewidth}\raggedright
Signif.
\end{minipage} \\
\midrule\noalign{}
\endhead
\bottomrule\noalign{}
\endlastfoot
Intercept & -4.258 & 0.547 & -7.782 & \textless0.001 & \\
np.log(capital) & 0.140 & 0.030 & 4.736 & \textless0.001 & *** \\
np.log(labour) & 0.473 & 0.197 & 2.404 & 0.017 & ** \\
np.log(wage) & 0.853 & 0.086 & 9.929 & \textless0.001 & *** \\
I(np.log(capital) ** 2) & 0.008 & 0.008 & 0.932 & 0.352 & \\
I(np.log(labour) ** 2) & 0.027 & 0.020 & 1.376 & 0.169 & \\
\end{longtable}
}

\emph{Signif. codes:} \texttt{***} \(p<0.01\), \texttt{**} \(p<0.05\),
\texttt{*} \(p<0.1\)

Модель была подогнана по 569 наблюдениям.

\begin{tcolorbox}[enhanced jigsaw, arc=.35mm, bottomrule=.15mm, bottomtitle=1mm, breakable, colback=white, colbacktitle=quarto-callout-note-color!10!white, colframe=quarto-callout-note-color-frame, coltitle=black, left=2mm, leftrule=.75mm, opacityback=0, opacitybacktitle=0.6, rightrule=.15mm, title=\textcolor{quarto-callout-note-color}{\faInfo}\hspace{0.5em}{Критическое значение}, titlerule=0mm, toprule=.15mm, toptitle=1mm]

\(t_{crit} = t_{0.95}(563) = 1.648\)

\end{tcolorbox}

\begin{tcolorbox}[enhanced jigsaw, arc=.35mm, bottomrule=.15mm, bottomtitle=1mm, breakable, colback=white, colbacktitle=quarto-callout-note-color!10!white, colframe=quarto-callout-note-color-frame, coltitle=black, left=2mm, leftrule=.75mm, opacityback=0, opacitybacktitle=0.6, rightrule=.15mm, title=\textcolor{quarto-callout-note-color}{\faInfo}\hspace{0.5em}{Коэффициенты, значимые на уровне 10\% (по робастным стандартным ошибкам
HC3)}, titlerule=0mm, toprule=.15mm, toptitle=1mm]

\texttt{np.log(capital)}, \texttt{np.log(labour)}, \texttt{np.log(wage)}

\end{tcolorbox}

\subsection{cost equation \#1}\label{cost-equation-1}

Для набора данных \texttt{Electricity} рассмотрим регрессию

\begin{Shaded}
\begin{Highlighting}[]
\NormalTok{np.log(cost) }\OperatorTok{\textasciitilde{}}\NormalTok{ np.log(q) }\OperatorTok{+}\NormalTok{ I(np.log(q)}\OperatorTok{**}\DecValTok{2}\NormalTok{) }\OperatorTok{+}
\NormalTok{    np.log(pl) }\OperatorTok{+}\NormalTok{ np.log(pk) }\OperatorTok{+}\NormalTok{ np.log(pf)}
\end{Highlighting}
\end{Shaded}

Переменные: \texttt{cost} --- общие издержки, \texttt{q} --- выпуск,
\texttt{pl}, \texttt{pk}, \texttt{pf} --- цены труда, капитала и
топлива.

\begin{tcolorbox}[enhanced jigsaw, arc=.35mm, bottomrule=.15mm, bottomtitle=1mm, breakable, colback=white, colbacktitle=quarto-callout-important-color!10!white, colframe=quarto-callout-important-color-frame, coltitle=black, left=2mm, leftrule=.75mm, opacityback=0, opacitybacktitle=0.6, rightrule=.15mm, title=\textcolor{quarto-callout-important-color}{\faExclamation}\hspace{0.5em}{Условия теста}, titlerule=0mm, toprule=.15mm, toptitle=1mm]

\begin{itemize}
\tightlist
\item
  уровень значимости: \textbf{1\%}
\item
  стандартные ошибки: \textbf{неробастные}
\item
  гипотеза: \(H_0:\beta_j=0\) для каждого коэффициента
\end{itemize}

\end{tcolorbox}

\begin{itemize}
\tightlist
\item
  Оцените регрессию.
\item
  Выведите результаты t-теста.
\item
  Вычислите критическое значение.
\item
  Какие коэффициенты значимы?
\end{itemize}

\textbf{Ответ округлите до 3-х десятичных знаков.}

\textbf{Результаты \(t\)-теста для коэффициентов (неробастные
стандартные ошибки):}

{\def\LTcaptype{none} % do not increment counter
\begin{longtable}[]{@{}
  >{\raggedright\arraybackslash}p{(\linewidth - 10\tabcolsep) * \real{0.2603}}
  >{\raggedleft\arraybackslash}p{(\linewidth - 10\tabcolsep) * \real{0.1233}}
  >{\raggedleft\arraybackslash}p{(\linewidth - 10\tabcolsep) * \real{0.1644}}
  >{\raggedleft\arraybackslash}p{(\linewidth - 10\tabcolsep) * \real{0.1507}}
  >{\raggedleft\arraybackslash}p{(\linewidth - 10\tabcolsep) * \real{0.1507}}
  >{\raggedright\arraybackslash}p{(\linewidth - 10\tabcolsep) * \real{0.1507}}@{}}
\toprule\noalign{}
\begin{minipage}[b]{\linewidth}\raggedright
\end{minipage} & \begin{minipage}[b]{\linewidth}\raggedleft
Coef.
\end{minipage} & \begin{minipage}[b]{\linewidth}\raggedleft
Std.Err.
\end{minipage} & \begin{minipage}[b]{\linewidth}\raggedleft
t-value
\end{minipage} & \begin{minipage}[b]{\linewidth}\raggedleft
p-value
\end{minipage} & \begin{minipage}[b]{\linewidth}\raggedright
Signif.
\end{minipage} \\
\midrule\noalign{}
\endhead
\bottomrule\noalign{}
\endlastfoot
Intercept & -6.739 & 0.706 & -9.541 & \textless0.001 & \\
np.log(q) & 0.403 & 0.032 & 12.734 & \textless0.001 & *** \\
I(np.log(q) ** 2) & 0.030 & 0.002 & 14.024 & \textless0.001 & *** \\
np.log(pl) & 0.146 & 0.070 & 2.073 & 0.040 & ** \\
np.log(pk) & 0.157 & 0.058 & 2.721 & 0.007 & *** \\
np.log(pf) & 0.685 & 0.043 & 16.043 & \textless0.001 & *** \\
\end{longtable}
}

\emph{Signif. codes:} \texttt{***} \(p<0.01\), \texttt{**} \(p<0.05\),
\texttt{*} \(p<0.1\)

Модель была подогнана по 158 наблюдениям.

\begin{tcolorbox}[enhanced jigsaw, arc=.35mm, bottomrule=.15mm, bottomtitle=1mm, breakable, colback=white, colbacktitle=quarto-callout-note-color!10!white, colframe=quarto-callout-note-color-frame, coltitle=black, left=2mm, leftrule=.75mm, opacityback=0, opacitybacktitle=0.6, rightrule=.15mm, title=\textcolor{quarto-callout-note-color}{\faInfo}\hspace{0.5em}{Критическое значение}, titlerule=0mm, toprule=.15mm, toptitle=1mm]

\(t_{crit} = t_{0.995}(152) = 2.609\)

\end{tcolorbox}

\begin{tcolorbox}[enhanced jigsaw, arc=.35mm, bottomrule=.15mm, bottomtitle=1mm, breakable, colback=white, colbacktitle=quarto-callout-note-color!10!white, colframe=quarto-callout-note-color-frame, coltitle=black, left=2mm, leftrule=.75mm, opacityback=0, opacitybacktitle=0.6, rightrule=.15mm, title=\textcolor{quarto-callout-note-color}{\faInfo}\hspace{0.5em}{Коэффициенты, значимые на уровне 1\% (по неробастным стандартным
ошибкам)}, titlerule=0mm, toprule=.15mm, toptitle=1mm]

\texttt{np.log(q)}, \texttt{I(np.log(q)\ **\ 2)}, \texttt{np.log(pk)},
\texttt{np.log(pf)}

\end{tcolorbox}

\subsection{cost equation \#2}\label{cost-equation-2}

Для набора данных \texttt{Electricity} рассмотрим регрессию

\begin{Shaded}
\begin{Highlighting}[]
\NormalTok{np.log(cost) }\OperatorTok{\textasciitilde{}}\NormalTok{ np.log(q) }\OperatorTok{+}\NormalTok{ I(np.log(q)}\OperatorTok{**}\DecValTok{2}\NormalTok{) }\OperatorTok{+}
\NormalTok{    np.log(pl) }\OperatorTok{+}\NormalTok{ np.log(pk) }\OperatorTok{+}\NormalTok{ np.log(pf)}
\end{Highlighting}
\end{Shaded}

\begin{tcolorbox}[enhanced jigsaw, arc=.35mm, bottomrule=.15mm, bottomtitle=1mm, breakable, colback=white, colbacktitle=quarto-callout-important-color!10!white, colframe=quarto-callout-important-color-frame, coltitle=black, left=2mm, leftrule=.75mm, opacityback=0, opacitybacktitle=0.6, rightrule=.15mm, title=\textcolor{quarto-callout-important-color}{\faExclamation}\hspace{0.5em}{Условия теста}, titlerule=0mm, toprule=.15mm, toptitle=1mm]

\begin{itemize}
\tightlist
\item
  уровень значимости: \textbf{1\%}
\item
  стандартные ошибки: \textbf{робастные (HC3)}
\item
  гипотеза: \(H_0:\beta_j=0\) для каждого коэффициента
\end{itemize}

\end{tcolorbox}

\begin{itemize}
\tightlist
\item
  Оцените регрессию.
\item
  Выведите результаты t-теста.
\item
  Вычислите критическое значение.
\item
  Какие коэффициенты значимы?
\end{itemize}

\textbf{Ответ округлите до 3-х десятичных знаков.}

\textbf{Результаты \(t\)-теста для коэффициентов (робастные стандартные
ошибки, HC3):}

{\def\LTcaptype{none} % do not increment counter
\begin{longtable}[]{@{}
  >{\raggedright\arraybackslash}p{(\linewidth - 10\tabcolsep) * \real{0.2603}}
  >{\raggedleft\arraybackslash}p{(\linewidth - 10\tabcolsep) * \real{0.1233}}
  >{\raggedleft\arraybackslash}p{(\linewidth - 10\tabcolsep) * \real{0.1644}}
  >{\raggedleft\arraybackslash}p{(\linewidth - 10\tabcolsep) * \real{0.1507}}
  >{\raggedleft\arraybackslash}p{(\linewidth - 10\tabcolsep) * \real{0.1507}}
  >{\raggedright\arraybackslash}p{(\linewidth - 10\tabcolsep) * \real{0.1507}}@{}}
\toprule\noalign{}
\begin{minipage}[b]{\linewidth}\raggedright
\end{minipage} & \begin{minipage}[b]{\linewidth}\raggedleft
Coef.
\end{minipage} & \begin{minipage}[b]{\linewidth}\raggedleft
Std.Err.
\end{minipage} & \begin{minipage}[b]{\linewidth}\raggedleft
t-value
\end{minipage} & \begin{minipage}[b]{\linewidth}\raggedleft
p-value
\end{minipage} & \begin{minipage}[b]{\linewidth}\raggedright
Signif.
\end{minipage} \\
\midrule\noalign{}
\endhead
\bottomrule\noalign{}
\endlastfoot
Intercept & -6.739 & 0.847 & -7.954 & \textless0.001 & \\
np.log(q) & 0.403 & 0.066 & 6.086 & \textless0.001 & *** \\
I(np.log(q) ** 2) & 0.030 & 0.004 & 7.419 & \textless0.001 & *** \\
np.log(pl) & 0.146 & 0.085 & 1.711 & 0.089 & * \\
np.log(pk) & 0.157 & 0.062 & 2.522 & 0.013 & ** \\
np.log(pf) & 0.685 & 0.052 & 13.188 & \textless0.001 & *** \\
\end{longtable}
}

\emph{Signif. codes:} \texttt{***} \(p<0.01\), \texttt{**} \(p<0.05\),
\texttt{*} \(p<0.1\)

Модель была подогнана по 158 наблюдениям.

\begin{tcolorbox}[enhanced jigsaw, arc=.35mm, bottomrule=.15mm, bottomtitle=1mm, breakable, colback=white, colbacktitle=quarto-callout-note-color!10!white, colframe=quarto-callout-note-color-frame, coltitle=black, left=2mm, leftrule=.75mm, opacityback=0, opacitybacktitle=0.6, rightrule=.15mm, title=\textcolor{quarto-callout-note-color}{\faInfo}\hspace{0.5em}{Критическое значение}, titlerule=0mm, toprule=.15mm, toptitle=1mm]

\(t_{crit} = t_{0.995}(152) = 2.609\)

\end{tcolorbox}

\begin{tcolorbox}[enhanced jigsaw, arc=.35mm, bottomrule=.15mm, bottomtitle=1mm, breakable, colback=white, colbacktitle=quarto-callout-note-color!10!white, colframe=quarto-callout-note-color-frame, coltitle=black, left=2mm, leftrule=.75mm, opacityback=0, opacitybacktitle=0.6, rightrule=.15mm, title=\textcolor{quarto-callout-note-color}{\faInfo}\hspace{0.5em}{Коэффициенты, значимые на уровне 1\% (по робастным стандартным ошибкам
HC3)}, titlerule=0mm, toprule=.15mm, toptitle=1mm]

\texttt{np.log(q)}, \texttt{I(np.log(q)\ **\ 2)}, \texttt{np.log(pf)}

\end{tcolorbox}

\section{t-тест: равенство коэффициента заданному
числу}\label{t-ux442ux435ux441ux442-ux440ux430ux432ux435ux43dux441ux442ux432ux43e-ux43aux43eux44dux444ux444ux438ux446ux438ux435ux43dux442ux430-ux437ux430ux434ux430ux43dux43dux43eux43cux443-ux447ux438ux441ux43bux443}

\subsection{output equation \#3}\label{output-equation-3}

Для набора данных \texttt{Labour} рассмотрим регрессию
\textbf{np.log(output) на np.log(capital), np.log(labour)}.

Результаты оценивания:

{\def\LTcaptype{none} % do not increment counter
\begin{longtable}[]{@{}lll@{}}
\toprule\noalign{}
& (1) & (2) \\
\midrule\noalign{}
\endhead
\bottomrule\noalign{}
\endlastfoot
Dep. Variable & np.log(output) & np.log(output) \\
np.log(capital) & 0.2076*** & 0.2076*** \\
& (0.01719) & (0.03159) \\
np.log(labour) & 0.7148*** & 0.7148*** \\
& (0.02314) & (0.04362) \\
Intercept & -1.711*** & -1.711*** \\
& (0.09671) & (0.1841) \\
R-squared & 0.8378 & 0.8378 \\
R-squared Adj. & 0.8373 & 0.8373 \\
F Statistic & 1462 & 486.1 \\
Residual Std. Error & 0.4781 & 0.4781 \\
\end{longtable}
}

\emph{Signif. codes:} \texttt{***} \(p<0.01\), \texttt{**} \(p<0.05\),
\texttt{*} \(p<0.1\). В скобках --- стандартные ошибки: (1) неробастные,
(2) HC3.

Модель была подогнана по 569 наблюдениям.

\begin{tcolorbox}[enhanced jigsaw, arc=.35mm, bottomrule=.15mm, bottomtitle=1mm, breakable, colback=white, colbacktitle=quarto-callout-important-color!10!white, colframe=quarto-callout-important-color-frame, coltitle=black, left=2mm, leftrule=.75mm, opacityback=0, opacitybacktitle=0.6, rightrule=.15mm, title=\textcolor{quarto-callout-important-color}{\faExclamation}\hspace{0.5em}{Условия теста}, titlerule=0mm, toprule=.15mm, toptitle=1mm]

\begin{itemize}
\tightlist
\item
  уровень значимости: \textbf{1\%}
\item
  стандартные ошибки: \textbf{неробастные и робастные (HC3)}
\item
  гипотеза: \(H_0:\beta_{capital}=0.5\)
\end{itemize}

\end{tcolorbox}

Вычислите тестовую статистику. Вычислите критическое значение. Какой
можно сделать вывод?

\textbf{Ответ округлите до 2-х десятичных знаков.}

\begin{tcolorbox}[enhanced jigsaw, arc=.35mm, bottomrule=.15mm, bottomtitle=1mm, breakable, colback=white, colbacktitle=quarto-callout-note-color!10!white, colframe=quarto-callout-note-color-frame, coltitle=black, left=2mm, leftrule=.75mm, opacityback=0, opacitybacktitle=0.6, rightrule=.15mm, title=\textcolor{quarto-callout-note-color}{\faInfo}\hspace{0.5em}{Критическое значение}, titlerule=0mm, toprule=.15mm, toptitle=1mm]

\(t_{crit} = t_{0.995}(566) = 2.58\)

\end{tcolorbox}

\textbf{Неробастный t-тест}

{\def\LTcaptype{none} % do not increment counter
\begin{longtable}[]{@{}
  >{\raggedright\arraybackslash}p{(\linewidth - 6\tabcolsep) * \real{0.2361}}
  >{\raggedleft\arraybackslash}p{(\linewidth - 6\tabcolsep) * \real{0.1389}}
  >{\raggedleft\arraybackslash}p{(\linewidth - 6\tabcolsep) * \real{0.3056}}
  >{\raggedleft\arraybackslash}p{(\linewidth - 6\tabcolsep) * \real{0.3194}}@{}}
\toprule\noalign{}
\begin{minipage}[b]{\linewidth}\raggedright
Переменная
\end{minipage} & \begin{minipage}[b]{\linewidth}\raggedleft
Оценка
\end{minipage} & \begin{minipage}[b]{\linewidth}\raggedleft
Стандартная ошибка
\end{minipage} & \begin{minipage}[b]{\linewidth}\raggedleft
Тестовая статистика
\end{minipage} \\
\midrule\noalign{}
\endhead
\bottomrule\noalign{}
\endlastfoot
np.log(capital) & 0.2076 & 0.01719 & -17.01 \\
\end{longtable}
}

\begin{tcolorbox}[enhanced jigsaw, arc=.35mm, bottomrule=.15mm, bottomtitle=1mm, breakable, colback=white, colbacktitle=quarto-callout-tip-color!10!white, colframe=quarto-callout-tip-color-frame, coltitle=black, left=2mm, leftrule=.75mm, opacityback=0, opacitybacktitle=0.6, rightrule=.15mm, title=\textcolor{quarto-callout-tip-color}{\faLightbulb}\hspace{0.5em}{Вывод}, titlerule=0mm, toprule=.15mm, toptitle=1mm]

\(\lvert t\rvert = 17.01 > 2.58\), гипотеза \(H_0\) отвергается.

\end{tcolorbox}

\textbf{Робастный t-тест (HC3)}

{\def\LTcaptype{none} % do not increment counter
\begin{longtable}[]{@{}
  >{\raggedright\arraybackslash}p{(\linewidth - 6\tabcolsep) * \real{0.2179}}
  >{\raggedleft\arraybackslash}p{(\linewidth - 6\tabcolsep) * \real{0.1282}}
  >{\raggedleft\arraybackslash}p{(\linewidth - 6\tabcolsep) * \real{0.3590}}
  >{\raggedleft\arraybackslash}p{(\linewidth - 6\tabcolsep) * \real{0.2949}}@{}}
\toprule\noalign{}
\begin{minipage}[b]{\linewidth}\raggedright
Переменная
\end{minipage} & \begin{minipage}[b]{\linewidth}\raggedleft
Оценка
\end{minipage} & \begin{minipage}[b]{\linewidth}\raggedleft
Стандартная ошибка (HC3)
\end{minipage} & \begin{minipage}[b]{\linewidth}\raggedleft
Тестовая статистика
\end{minipage} \\
\midrule\noalign{}
\endhead
\bottomrule\noalign{}
\endlastfoot
np.log(capital) & 0.2076 & 0.03159 & -9.26 \\
\end{longtable}
}

\begin{tcolorbox}[enhanced jigsaw, arc=.35mm, bottomrule=.15mm, bottomtitle=1mm, breakable, colback=white, colbacktitle=quarto-callout-tip-color!10!white, colframe=quarto-callout-tip-color-frame, coltitle=black, left=2mm, leftrule=.75mm, opacityback=0, opacitybacktitle=0.6, rightrule=.15mm, title=\textcolor{quarto-callout-tip-color}{\faLightbulb}\hspace{0.5em}{Вывод}, titlerule=0mm, toprule=.15mm, toptitle=1mm]

\(\lvert t\rvert = 9.26 > 2.58\), гипотеза \(H_0\) отвергается.

\end{tcolorbox}

\subsection{output equation \#4}\label{output-equation-4}

Для набора данных \texttt{Labour} рассмотрим регрессию
\textbf{np.log(output) на np.log(capital), np.log(labour)}.

Результаты оценивания:

{\def\LTcaptype{none} % do not increment counter
\begin{longtable}[]{@{}lll@{}}
\toprule\noalign{}
& (1) & (2) \\
\midrule\noalign{}
\endhead
\bottomrule\noalign{}
\endlastfoot
Dep. Variable & np.log(output) & np.log(output) \\
np.log(capital) & 0.2076*** & 0.2076*** \\
& (0.01719) & (0.03159) \\
np.log(labour) & 0.7148*** & 0.7148*** \\
& (0.02314) & (0.04362) \\
Intercept & -1.711*** & -1.711*** \\
& (0.09671) & (0.1841) \\
R-squared & 0.8378 & 0.8378 \\
R-squared Adj. & 0.8373 & 0.8373 \\
F Statistic & 1462 & 486.1 \\
Residual Std. Error & 0.4781 & 0.4781 \\
\end{longtable}
}

\emph{Signif. codes:} \texttt{***} \(p<0.01\), \texttt{**} \(p<0.05\),
\texttt{*} \(p<0.1\). В скобках --- стандартные ошибки: (1) неробастные,
(2) HC3.

Модель была подогнана по 569 наблюдениям.

\begin{tcolorbox}[enhanced jigsaw, arc=.35mm, bottomrule=.15mm, bottomtitle=1mm, breakable, colback=white, colbacktitle=quarto-callout-important-color!10!white, colframe=quarto-callout-important-color-frame, coltitle=black, left=2mm, leftrule=.75mm, opacityback=0, opacitybacktitle=0.6, rightrule=.15mm, title=\textcolor{quarto-callout-important-color}{\faExclamation}\hspace{0.5em}{Условия теста}, titlerule=0mm, toprule=.15mm, toptitle=1mm]

\begin{itemize}
\tightlist
\item
  уровень значимости: \textbf{10\%}
\item
  стандартные ошибки: \textbf{неробастные и робастные (HC3)}
\item
  гипотеза: \(H_0:\beta_{labour}=0.7\)
\end{itemize}

\end{tcolorbox}

Вычислите тестовую статистику. Вычислите критическое значение. Какой
можно сделать вывод?

\textbf{Ответ округлите до 2-х десятичных знаков.}

\begin{tcolorbox}[enhanced jigsaw, arc=.35mm, bottomrule=.15mm, bottomtitle=1mm, breakable, colback=white, colbacktitle=quarto-callout-note-color!10!white, colframe=quarto-callout-note-color-frame, coltitle=black, left=2mm, leftrule=.75mm, opacityback=0, opacitybacktitle=0.6, rightrule=.15mm, title=\textcolor{quarto-callout-note-color}{\faInfo}\hspace{0.5em}{Критическое значение}, titlerule=0mm, toprule=.15mm, toptitle=1mm]

\(t_{crit} = t_{0.95}(566) = 1.65\)

\end{tcolorbox}

\textbf{Неробастный t-тест}

{\def\LTcaptype{none} % do not increment counter
\begin{longtable}[]{@{}
  >{\raggedright\arraybackslash}p{(\linewidth - 6\tabcolsep) * \real{0.2254}}
  >{\raggedleft\arraybackslash}p{(\linewidth - 6\tabcolsep) * \real{0.1408}}
  >{\raggedleft\arraybackslash}p{(\linewidth - 6\tabcolsep) * \real{0.3099}}
  >{\raggedleft\arraybackslash}p{(\linewidth - 6\tabcolsep) * \real{0.3239}}@{}}
\toprule\noalign{}
\begin{minipage}[b]{\linewidth}\raggedright
Переменная
\end{minipage} & \begin{minipage}[b]{\linewidth}\raggedleft
Оценка
\end{minipage} & \begin{minipage}[b]{\linewidth}\raggedleft
Стандартная ошибка
\end{minipage} & \begin{minipage}[b]{\linewidth}\raggedleft
Тестовая статистика
\end{minipage} \\
\midrule\noalign{}
\endhead
\bottomrule\noalign{}
\endlastfoot
np.log(labour) & 0.7148 & 0.02314 & 0.64 \\
\end{longtable}
}

\begin{tcolorbox}[enhanced jigsaw, arc=.35mm, bottomrule=.15mm, bottomtitle=1mm, breakable, colback=white, colbacktitle=quarto-callout-tip-color!10!white, colframe=quarto-callout-tip-color-frame, coltitle=black, left=2mm, leftrule=.75mm, opacityback=0, opacitybacktitle=0.6, rightrule=.15mm, title=\textcolor{quarto-callout-tip-color}{\faLightbulb}\hspace{0.5em}{Вывод}, titlerule=0mm, toprule=.15mm, toptitle=1mm]

\(\lvert t\rvert = 0.64 < 1.65\), гипотеза \(H_0\) не отвергается.

\end{tcolorbox}

\textbf{Робастный t-тест (HC3)}

{\def\LTcaptype{none} % do not increment counter
\begin{longtable}[]{@{}
  >{\raggedright\arraybackslash}p{(\linewidth - 6\tabcolsep) * \real{0.2078}}
  >{\raggedleft\arraybackslash}p{(\linewidth - 6\tabcolsep) * \real{0.1299}}
  >{\raggedleft\arraybackslash}p{(\linewidth - 6\tabcolsep) * \real{0.3636}}
  >{\raggedleft\arraybackslash}p{(\linewidth - 6\tabcolsep) * \real{0.2987}}@{}}
\toprule\noalign{}
\begin{minipage}[b]{\linewidth}\raggedright
Переменная
\end{minipage} & \begin{minipage}[b]{\linewidth}\raggedleft
Оценка
\end{minipage} & \begin{minipage}[b]{\linewidth}\raggedleft
Стандартная ошибка (HC3)
\end{minipage} & \begin{minipage}[b]{\linewidth}\raggedleft
Тестовая статистика
\end{minipage} \\
\midrule\noalign{}
\endhead
\bottomrule\noalign{}
\endlastfoot
np.log(labour) & 0.7148 & 0.04362 & 0.34 \\
\end{longtable}
}

\begin{tcolorbox}[enhanced jigsaw, arc=.35mm, bottomrule=.15mm, bottomtitle=1mm, breakable, colback=white, colbacktitle=quarto-callout-tip-color!10!white, colframe=quarto-callout-tip-color-frame, coltitle=black, left=2mm, leftrule=.75mm, opacityback=0, opacitybacktitle=0.6, rightrule=.15mm, title=\textcolor{quarto-callout-tip-color}{\faLightbulb}\hspace{0.5em}{Вывод}, titlerule=0mm, toprule=.15mm, toptitle=1mm]

\(\lvert t\rvert = 0.34 < 1.65\), гипотеза \(H_0\) не отвергается.

\end{tcolorbox}

\subsection{cost equation \#3}\label{cost-equation-3}

Для набора данных \texttt{Electricity} рассмотрим регрессию
\textbf{np.log(cost) на np.log(q), np.log(pl), np.log(pk), np.log(pf)}.

Результаты оценивания:

{\def\LTcaptype{none} % do not increment counter
\begin{longtable}[]{@{}lll@{}}
\toprule\noalign{}
& (1) & (2) \\
\midrule\noalign{}
\endhead
\bottomrule\noalign{}
\endlastfoot
Dep. Variable & np.log(cost) & np.log(cost) \\
np.log(q) & 0.8381*** & 0.8381*** \\
& (0.009387) & (0.0205) \\
np.log(pl) & 0.04426 & 0.04426 \\
& (0.1058) & (0.119) \\
np.log(pk) & 0.1884** & 0.1884** \\
& (0.08707) & (0.089) \\
np.log(pf) & 0.7131*** & 0.7131*** \\
& (0.06436) & (0.0814) \\
Intercept & -7.472*** & -7.472*** \\
& (1.063) & (1.156) \\
R-squared & 0.9822 & 0.9822 \\
R-squared Adj. & 0.9817 & 0.9817 \\
F Statistic & 2107 & 552.7 \\
Residual Std. Error & 0.2087 & 0.2087 \\
\end{longtable}
}

\emph{Signif. codes:} \texttt{***} \(p<0.01\), \texttt{**} \(p<0.05\),
\texttt{*} \(p<0.1\). В скобках --- стандартные ошибки: (1) неробастные,
(2) HC3.

Модель была подогнана по 158 наблюдениям.

\begin{tcolorbox}[enhanced jigsaw, arc=.35mm, bottomrule=.15mm, bottomtitle=1mm, breakable, colback=white, colbacktitle=quarto-callout-important-color!10!white, colframe=quarto-callout-important-color-frame, coltitle=black, left=2mm, leftrule=.75mm, opacityback=0, opacitybacktitle=0.6, rightrule=.15mm, title=\textcolor{quarto-callout-important-color}{\faExclamation}\hspace{0.5em}{Условия теста}, titlerule=0mm, toprule=.15mm, toptitle=1mm]

\begin{itemize}
\tightlist
\item
  уровень значимости: \textbf{5\%}
\item
  стандартные ошибки: \textbf{неробастные и робастные (HC3)}
\item
  гипотеза: \(H_0:\beta_{q}=1\)
\end{itemize}

\end{tcolorbox}

Вычислите тестовую статистику. Вычислите критическое значение. Какой
можно сделать вывод?

\textbf{Ответ округлите до 2-х десятичных знаков.}

\begin{tcolorbox}[enhanced jigsaw, arc=.35mm, bottomrule=.15mm, bottomtitle=1mm, breakable, colback=white, colbacktitle=quarto-callout-note-color!10!white, colframe=quarto-callout-note-color-frame, coltitle=black, left=2mm, leftrule=.75mm, opacityback=0, opacitybacktitle=0.6, rightrule=.15mm, title=\textcolor{quarto-callout-note-color}{\faInfo}\hspace{0.5em}{Критическое значение}, titlerule=0mm, toprule=.15mm, toptitle=1mm]

\(t_{crit} = t_{0.975}(153) = 1.98\)

\end{tcolorbox}

\textbf{Неробастный t-тест}

{\def\LTcaptype{none} % do not increment counter
\begin{longtable}[]{@{}
  >{\raggedright\arraybackslash}p{(\linewidth - 6\tabcolsep) * \real{0.2029}}
  >{\raggedleft\arraybackslash}p{(\linewidth - 6\tabcolsep) * \real{0.1449}}
  >{\raggedleft\arraybackslash}p{(\linewidth - 6\tabcolsep) * \real{0.3188}}
  >{\raggedleft\arraybackslash}p{(\linewidth - 6\tabcolsep) * \real{0.3333}}@{}}
\toprule\noalign{}
\begin{minipage}[b]{\linewidth}\raggedright
Переменная
\end{minipage} & \begin{minipage}[b]{\linewidth}\raggedleft
Оценка
\end{minipage} & \begin{minipage}[b]{\linewidth}\raggedleft
Стандартная ошибка
\end{minipage} & \begin{minipage}[b]{\linewidth}\raggedleft
Тестовая статистика
\end{minipage} \\
\midrule\noalign{}
\endhead
\bottomrule\noalign{}
\endlastfoot
np.log(q) & 0.8381 & 0.009387 & -17.25 \\
\end{longtable}
}

\begin{tcolorbox}[enhanced jigsaw, arc=.35mm, bottomrule=.15mm, bottomtitle=1mm, breakable, colback=white, colbacktitle=quarto-callout-tip-color!10!white, colframe=quarto-callout-tip-color-frame, coltitle=black, left=2mm, leftrule=.75mm, opacityback=0, opacitybacktitle=0.6, rightrule=.15mm, title=\textcolor{quarto-callout-tip-color}{\faLightbulb}\hspace{0.5em}{Вывод}, titlerule=0mm, toprule=.15mm, toptitle=1mm]

\(\lvert t\rvert = 17.25 > 1.98\), гипотеза \(H_0\) отвергается.

\end{tcolorbox}

\textbf{Робастный t-тест (HC3)}

{\def\LTcaptype{none} % do not increment counter
\begin{longtable}[]{@{}
  >{\raggedright\arraybackslash}p{(\linewidth - 6\tabcolsep) * \real{0.1867}}
  >{\raggedleft\arraybackslash}p{(\linewidth - 6\tabcolsep) * \real{0.1333}}
  >{\raggedleft\arraybackslash}p{(\linewidth - 6\tabcolsep) * \real{0.3733}}
  >{\raggedleft\arraybackslash}p{(\linewidth - 6\tabcolsep) * \real{0.3067}}@{}}
\toprule\noalign{}
\begin{minipage}[b]{\linewidth}\raggedright
Переменная
\end{minipage} & \begin{minipage}[b]{\linewidth}\raggedleft
Оценка
\end{minipage} & \begin{minipage}[b]{\linewidth}\raggedleft
Стандартная ошибка (HC3)
\end{minipage} & \begin{minipage}[b]{\linewidth}\raggedleft
Тестовая статистика
\end{minipage} \\
\midrule\noalign{}
\endhead
\bottomrule\noalign{}
\endlastfoot
np.log(q) & 0.8381 & 0.0205 & -7.90 \\
\end{longtable}
}

\begin{tcolorbox}[enhanced jigsaw, arc=.35mm, bottomrule=.15mm, bottomtitle=1mm, breakable, colback=white, colbacktitle=quarto-callout-tip-color!10!white, colframe=quarto-callout-tip-color-frame, coltitle=black, left=2mm, leftrule=.75mm, opacityback=0, opacitybacktitle=0.6, rightrule=.15mm, title=\textcolor{quarto-callout-tip-color}{\faLightbulb}\hspace{0.5em}{Вывод}, titlerule=0mm, toprule=.15mm, toptitle=1mm]

\(\lvert t\rvert = 7.90 > 1.98\), гипотеза \(H_0\) отвергается.

\end{tcolorbox}

\subsection{cost equation \#4}\label{cost-equation-4}

Для набора данных \texttt{Electricity} рассмотрим регрессию
\textbf{np.log(cost) на np.log(q), np.log(pl), np.log(pk), np.log(pf)}.

Результаты оценивания:

{\def\LTcaptype{none} % do not increment counter
\begin{longtable}[]{@{}lll@{}}
\toprule\noalign{}
& (1) & (2) \\
\midrule\noalign{}
\endhead
\bottomrule\noalign{}
\endlastfoot
Dep. Variable & np.log(cost) & np.log(cost) \\
np.log(q) & 0.8381*** & 0.8381*** \\
& (0.009387) & (0.0205) \\
np.log(pl) & 0.04426 & 0.04426 \\
& (0.1058) & (0.119) \\
np.log(pk) & 0.1884** & 0.1884** \\
& (0.08707) & (0.089) \\
np.log(pf) & 0.7131*** & 0.7131*** \\
& (0.06436) & (0.0814) \\
Intercept & -7.472*** & -7.472*** \\
& (1.063) & (1.156) \\
R-squared & 0.9822 & 0.9822 \\
R-squared Adj. & 0.9817 & 0.9817 \\
F Statistic & 2107 & 552.7 \\
Residual Std. Error & 0.2087 & 0.2087 \\
\end{longtable}
}

\emph{Signif. codes:} \texttt{***} \(p<0.01\), \texttt{**} \(p<0.05\),
\texttt{*} \(p<0.1\). В скобках --- стандартные ошибки: (1) неробастные,
(2) HC3.

Модель была подогнана по 158 наблюдениям.

\begin{tcolorbox}[enhanced jigsaw, arc=.35mm, bottomrule=.15mm, bottomtitle=1mm, breakable, colback=white, colbacktitle=quarto-callout-important-color!10!white, colframe=quarto-callout-important-color-frame, coltitle=black, left=2mm, leftrule=.75mm, opacityback=0, opacitybacktitle=0.6, rightrule=.15mm, title=\textcolor{quarto-callout-important-color}{\faExclamation}\hspace{0.5em}{Условия теста}, titlerule=0mm, toprule=.15mm, toptitle=1mm]

\begin{itemize}
\tightlist
\item
  уровень значимости: \textbf{10\%}
\item
  стандартные ошибки: \textbf{неробастные и робастные (HC3)}
\item
  гипотеза: \(H_0:\beta_{pf}=1\)
\end{itemize}

\end{tcolorbox}

Вычислите тестовую статистику. Вычислите критическое значение. Какой
можно сделать вывод?

\textbf{Ответ округлите до 2-х десятичных знаков.}

\begin{tcolorbox}[enhanced jigsaw, arc=.35mm, bottomrule=.15mm, bottomtitle=1mm, breakable, colback=white, colbacktitle=quarto-callout-note-color!10!white, colframe=quarto-callout-note-color-frame, coltitle=black, left=2mm, leftrule=.75mm, opacityback=0, opacitybacktitle=0.6, rightrule=.15mm, title=\textcolor{quarto-callout-note-color}{\faInfo}\hspace{0.5em}{Критическое значение}, titlerule=0mm, toprule=.15mm, toptitle=1mm]

\(t_{crit} = t_{0.95}(153) = 1.65\)

\end{tcolorbox}

\textbf{Неробастный t-тест}

{\def\LTcaptype{none} % do not increment counter
\begin{longtable}[]{@{}
  >{\raggedright\arraybackslash}p{(\linewidth - 6\tabcolsep) * \real{0.2029}}
  >{\raggedleft\arraybackslash}p{(\linewidth - 6\tabcolsep) * \real{0.1449}}
  >{\raggedleft\arraybackslash}p{(\linewidth - 6\tabcolsep) * \real{0.3188}}
  >{\raggedleft\arraybackslash}p{(\linewidth - 6\tabcolsep) * \real{0.3333}}@{}}
\toprule\noalign{}
\begin{minipage}[b]{\linewidth}\raggedright
Переменная
\end{minipage} & \begin{minipage}[b]{\linewidth}\raggedleft
Оценка
\end{minipage} & \begin{minipage}[b]{\linewidth}\raggedleft
Стандартная ошибка
\end{minipage} & \begin{minipage}[b]{\linewidth}\raggedleft
Тестовая статистика
\end{minipage} \\
\midrule\noalign{}
\endhead
\bottomrule\noalign{}
\endlastfoot
np.log(pf) & 0.7131 & 0.06436 & -4.46 \\
\end{longtable}
}

\begin{tcolorbox}[enhanced jigsaw, arc=.35mm, bottomrule=.15mm, bottomtitle=1mm, breakable, colback=white, colbacktitle=quarto-callout-tip-color!10!white, colframe=quarto-callout-tip-color-frame, coltitle=black, left=2mm, leftrule=.75mm, opacityback=0, opacitybacktitle=0.6, rightrule=.15mm, title=\textcolor{quarto-callout-tip-color}{\faLightbulb}\hspace{0.5em}{Вывод}, titlerule=0mm, toprule=.15mm, toptitle=1mm]

\(\lvert t\rvert = 4.46 > 1.65\), гипотеза \(H_0\) отвергается.

\end{tcolorbox}

\textbf{Робастный t-тест (HC3)}

{\def\LTcaptype{none} % do not increment counter
\begin{longtable}[]{@{}
  >{\raggedright\arraybackslash}p{(\linewidth - 6\tabcolsep) * \real{0.1867}}
  >{\raggedleft\arraybackslash}p{(\linewidth - 6\tabcolsep) * \real{0.1333}}
  >{\raggedleft\arraybackslash}p{(\linewidth - 6\tabcolsep) * \real{0.3733}}
  >{\raggedleft\arraybackslash}p{(\linewidth - 6\tabcolsep) * \real{0.3067}}@{}}
\toprule\noalign{}
\begin{minipage}[b]{\linewidth}\raggedright
Переменная
\end{minipage} & \begin{minipage}[b]{\linewidth}\raggedleft
Оценка
\end{minipage} & \begin{minipage}[b]{\linewidth}\raggedleft
Стандартная ошибка (HC3)
\end{minipage} & \begin{minipage}[b]{\linewidth}\raggedleft
Тестовая статистика
\end{minipage} \\
\midrule\noalign{}
\endhead
\bottomrule\noalign{}
\endlastfoot
np.log(pf) & 0.7131 & 0.0814 & -3.52 \\
\end{longtable}
}

\begin{tcolorbox}[enhanced jigsaw, arc=.35mm, bottomrule=.15mm, bottomtitle=1mm, breakable, colback=white, colbacktitle=quarto-callout-tip-color!10!white, colframe=quarto-callout-tip-color-frame, coltitle=black, left=2mm, leftrule=.75mm, opacityback=0, opacitybacktitle=0.6, rightrule=.15mm, title=\textcolor{quarto-callout-tip-color}{\faLightbulb}\hspace{0.5em}{Вывод}, titlerule=0mm, toprule=.15mm, toptitle=1mm]

\(\lvert t\rvert = 3.52 > 1.65\), гипотеза \(H_0\) отвергается.

\end{tcolorbox}

\subsection{sleep equation \#5}\label{sleep-equation-5}

Для набора данных \texttt{sleep75} рассмотрим регрессию

\texttt{sleep\ \textasciitilde{}\ totwrk\ +\ age\ +\ south\ +\ male\ +\ smsa\ +\ yngkid\ +\ marr\ +\ union}

Результаты оценивания:

{\def\LTcaptype{none} % do not increment counter
\begin{longtable}[]{@{}lll@{}}
\toprule\noalign{}
& (1) & (2) \\
\midrule\noalign{}
\endhead
\bottomrule\noalign{}
\endlastfoot
Dep. Variable & sleep & sleep \\
totwrk & -0.1691*** & -0.1691*** \\
& (0.01807) & (0.02035) \\
age & 2.714* & 2.714* \\
& (1.472) & (1.451) \\
south & 102.3** & 102.3** \\
& (41.93) & (42.59) \\
male & 87.11** & 87.11** \\
& (35.17) & (35.96) \\
smsa & -54.19 & -54.19 \\
& (33.19) & (33.45) \\
yngkid & -13.05 & -13.05 \\
& (50.46) & (54.58) \\
marr & 31.36 & 31.36 \\
& (42.26) & (46.27) \\
union & 11.87 & 11.87 \\
& (38.19) & (37.53) \\
Intercept & 3447*** & 3447*** \\
& (81.84) & (91.67) \\
R-squared & 0.1314 & 0.1314 \\
R-squared Adj. & 0.1214 & 0.1214 \\
F Statistic & 13.18 & 10.44 \\
Residual Std. Error & 416.6 & 416.6 \\
\end{longtable}
}

\emph{Signif. codes:} \texttt{***} \(p<0.01\), \texttt{**} \(p<0.05\),
\texttt{*} \(p<0.1\). В скобках --- стандартные ошибки: (1) неробастные,
(2) HC3.

Модель была подогнана по 706 наблюдениям.

\begin{tcolorbox}[enhanced jigsaw, arc=.35mm, bottomrule=.15mm, bottomtitle=1mm, breakable, colback=white, colbacktitle=quarto-callout-important-color!10!white, colframe=quarto-callout-important-color-frame, coltitle=black, left=2mm, leftrule=.75mm, opacityback=0, opacitybacktitle=0.6, rightrule=.15mm, title=\textcolor{quarto-callout-important-color}{\faExclamation}\hspace{0.5em}{Условия теста}, titlerule=0mm, toprule=.15mm, toptitle=1mm]

\begin{itemize}
\tightlist
\item
  уровень значимости: \textbf{1\%}
\item
  стандартные ошибки: \textbf{неробастные и робастные (HC3)}
\item
  гипотеза: \(H_0:\beta_{male}=90\)
\end{itemize}

\end{tcolorbox}

Вычислите тестовую статистику. Вычислите критическое значение. Какой
можно сделать вывод?

\textbf{Ответ округлите до 2-х десятичных знаков.}

\begin{tcolorbox}[enhanced jigsaw, arc=.35mm, bottomrule=.15mm, bottomtitle=1mm, breakable, colback=white, colbacktitle=quarto-callout-note-color!10!white, colframe=quarto-callout-note-color-frame, coltitle=black, left=2mm, leftrule=.75mm, opacityback=0, opacitybacktitle=0.6, rightrule=.15mm, title=\textcolor{quarto-callout-note-color}{\faInfo}\hspace{0.5em}{Критическое значение}, titlerule=0mm, toprule=.15mm, toptitle=1mm]

\(t_{crit} = t_{0.995}(697) = 2.58\)

\end{tcolorbox}

\textbf{Неробастный t-тест}

{\def\LTcaptype{none} % do not increment counter
\begin{longtable}[]{@{}
  >{\raggedright\arraybackslash}p{(\linewidth - 6\tabcolsep) * \real{0.2029}}
  >{\raggedleft\arraybackslash}p{(\linewidth - 6\tabcolsep) * \real{0.1449}}
  >{\raggedleft\arraybackslash}p{(\linewidth - 6\tabcolsep) * \real{0.3188}}
  >{\raggedleft\arraybackslash}p{(\linewidth - 6\tabcolsep) * \real{0.3333}}@{}}
\toprule\noalign{}
\begin{minipage}[b]{\linewidth}\raggedright
Переменная
\end{minipage} & \begin{minipage}[b]{\linewidth}\raggedleft
Оценка
\end{minipage} & \begin{minipage}[b]{\linewidth}\raggedleft
Стандартная ошибка
\end{minipage} & \begin{minipage}[b]{\linewidth}\raggedleft
Тестовая статистика
\end{minipage} \\
\midrule\noalign{}
\endhead
\bottomrule\noalign{}
\endlastfoot
male & 87.11 & 35.17 & -0.08 \\
\end{longtable}
}

\begin{tcolorbox}[enhanced jigsaw, arc=.35mm, bottomrule=.15mm, bottomtitle=1mm, breakable, colback=white, colbacktitle=quarto-callout-tip-color!10!white, colframe=quarto-callout-tip-color-frame, coltitle=black, left=2mm, leftrule=.75mm, opacityback=0, opacitybacktitle=0.6, rightrule=.15mm, title=\textcolor{quarto-callout-tip-color}{\faLightbulb}\hspace{0.5em}{Вывод}, titlerule=0mm, toprule=.15mm, toptitle=1mm]

\(\lvert t\rvert = 0.08 < 2.58\), гипотеза \(H_0\) не отвергается.

\end{tcolorbox}

\textbf{Робастный t-тест (HC3)}

{\def\LTcaptype{none} % do not increment counter
\begin{longtable}[]{@{}
  >{\raggedright\arraybackslash}p{(\linewidth - 6\tabcolsep) * \real{0.1867}}
  >{\raggedleft\arraybackslash}p{(\linewidth - 6\tabcolsep) * \real{0.1333}}
  >{\raggedleft\arraybackslash}p{(\linewidth - 6\tabcolsep) * \real{0.3733}}
  >{\raggedleft\arraybackslash}p{(\linewidth - 6\tabcolsep) * \real{0.3067}}@{}}
\toprule\noalign{}
\begin{minipage}[b]{\linewidth}\raggedright
Переменная
\end{minipage} & \begin{minipage}[b]{\linewidth}\raggedleft
Оценка
\end{minipage} & \begin{minipage}[b]{\linewidth}\raggedleft
Стандартная ошибка (HC3)
\end{minipage} & \begin{minipage}[b]{\linewidth}\raggedleft
Тестовая статистика
\end{minipage} \\
\midrule\noalign{}
\endhead
\bottomrule\noalign{}
\endlastfoot
male & 87.11 & 35.96 & -0.08 \\
\end{longtable}
}

\begin{tcolorbox}[enhanced jigsaw, arc=.35mm, bottomrule=.15mm, bottomtitle=1mm, breakable, colback=white, colbacktitle=quarto-callout-tip-color!10!white, colframe=quarto-callout-tip-color-frame, coltitle=black, left=2mm, leftrule=.75mm, opacityback=0, opacitybacktitle=0.6, rightrule=.15mm, title=\textcolor{quarto-callout-tip-color}{\faLightbulb}\hspace{0.5em}{Вывод}, titlerule=0mm, toprule=.15mm, toptitle=1mm]

\(\lvert t\rvert = 0.08 < 2.58\), гипотеза \(H_0\) не отвергается.

\end{tcolorbox}

\subsection{sleep equation \#6}\label{sleep-equation-6}

Для набора данных \texttt{sleep75} рассмотрим регрессию

\texttt{sleep\ \textasciitilde{}\ totwrk\ +\ age\ +\ south\ +\ male\ +\ smsa\ +\ yngkid\ +\ marr\ +\ union}

Результаты оценивания:

{\def\LTcaptype{none} % do not increment counter
\begin{longtable}[]{@{}lll@{}}
\toprule\noalign{}
& (1) & (2) \\
\midrule\noalign{}
\endhead
\bottomrule\noalign{}
\endlastfoot
Dep. Variable & sleep & sleep \\
totwrk & -0.1691*** & -0.1691*** \\
& (0.01807) & (0.02035) \\
age & 2.714* & 2.714* \\
& (1.472) & (1.451) \\
south & 102.3** & 102.3** \\
& (41.93) & (42.59) \\
male & 87.11** & 87.11** \\
& (35.17) & (35.96) \\
smsa & -54.19 & -54.19 \\
& (33.19) & (33.45) \\
yngkid & -13.05 & -13.05 \\
& (50.46) & (54.58) \\
marr & 31.36 & 31.36 \\
& (42.26) & (46.27) \\
union & 11.87 & 11.87 \\
& (38.19) & (37.53) \\
Intercept & 3447*** & 3447*** \\
& (81.84) & (91.67) \\
R-squared & 0.1314 & 0.1314 \\
R-squared Adj. & 0.1214 & 0.1214 \\
F Statistic & 13.18 & 10.44 \\
Residual Std. Error & 416.6 & 416.6 \\
\end{longtable}
}

\emph{Signif. codes:} \texttt{***} \(p<0.01\), \texttt{**} \(p<0.05\),
\texttt{*} \(p<0.1\). В скобках --- стандартные ошибки: (1) неробастные,
(2) HC3.

Модель была подогнана по 706 наблюдениям.

\begin{tcolorbox}[enhanced jigsaw, arc=.35mm, bottomrule=.15mm, bottomtitle=1mm, breakable, colback=white, colbacktitle=quarto-callout-important-color!10!white, colframe=quarto-callout-important-color-frame, coltitle=black, left=2mm, leftrule=.75mm, opacityback=0, opacitybacktitle=0.6, rightrule=.15mm, title=\textcolor{quarto-callout-important-color}{\faExclamation}\hspace{0.5em}{Условия теста}, titlerule=0mm, toprule=.15mm, toptitle=1mm]

\begin{itemize}
\tightlist
\item
  уровень значимости: \textbf{1\%}
\item
  стандартные ошибки: \textbf{неробастные и робастные (HC3)}
\item
  гипотеза: \(H_0:\beta_{south}=100\)
\end{itemize}

\end{tcolorbox}

Вычислите тестовую статистику. Вычислите критическое значение. Какой
можно сделать вывод?

\textbf{Ответ округлите до 2-х десятичных знаков.}

\begin{tcolorbox}[enhanced jigsaw, arc=.35mm, bottomrule=.15mm, bottomtitle=1mm, breakable, colback=white, colbacktitle=quarto-callout-note-color!10!white, colframe=quarto-callout-note-color-frame, coltitle=black, left=2mm, leftrule=.75mm, opacityback=0, opacitybacktitle=0.6, rightrule=.15mm, title=\textcolor{quarto-callout-note-color}{\faInfo}\hspace{0.5em}{Критическое значение}, titlerule=0mm, toprule=.15mm, toptitle=1mm]

\(t_{crit} = t_{0.995}(697) = 2.58\)

\end{tcolorbox}

\textbf{Неробастный t-тест}

{\def\LTcaptype{none} % do not increment counter
\begin{longtable}[]{@{}
  >{\raggedright\arraybackslash}p{(\linewidth - 6\tabcolsep) * \real{0.2029}}
  >{\raggedleft\arraybackslash}p{(\linewidth - 6\tabcolsep) * \real{0.1449}}
  >{\raggedleft\arraybackslash}p{(\linewidth - 6\tabcolsep) * \real{0.3188}}
  >{\raggedleft\arraybackslash}p{(\linewidth - 6\tabcolsep) * \real{0.3333}}@{}}
\toprule\noalign{}
\begin{minipage}[b]{\linewidth}\raggedright
Переменная
\end{minipage} & \begin{minipage}[b]{\linewidth}\raggedleft
Оценка
\end{minipage} & \begin{minipage}[b]{\linewidth}\raggedleft
Стандартная ошибка
\end{minipage} & \begin{minipage}[b]{\linewidth}\raggedleft
Тестовая статистика
\end{minipage} \\
\midrule\noalign{}
\endhead
\bottomrule\noalign{}
\endlastfoot
south & 102.3 & 41.93 & 0.05 \\
\end{longtable}
}

\begin{tcolorbox}[enhanced jigsaw, arc=.35mm, bottomrule=.15mm, bottomtitle=1mm, breakable, colback=white, colbacktitle=quarto-callout-tip-color!10!white, colframe=quarto-callout-tip-color-frame, coltitle=black, left=2mm, leftrule=.75mm, opacityback=0, opacitybacktitle=0.6, rightrule=.15mm, title=\textcolor{quarto-callout-tip-color}{\faLightbulb}\hspace{0.5em}{Вывод}, titlerule=0mm, toprule=.15mm, toptitle=1mm]

\(\lvert t\rvert = 0.05 < 2.58\), гипотеза \(H_0\) не отвергается.

\end{tcolorbox}

\textbf{Робастный t-тест (HC3)}

{\def\LTcaptype{none} % do not increment counter
\begin{longtable}[]{@{}
  >{\raggedright\arraybackslash}p{(\linewidth - 6\tabcolsep) * \real{0.1867}}
  >{\raggedleft\arraybackslash}p{(\linewidth - 6\tabcolsep) * \real{0.1333}}
  >{\raggedleft\arraybackslash}p{(\linewidth - 6\tabcolsep) * \real{0.3733}}
  >{\raggedleft\arraybackslash}p{(\linewidth - 6\tabcolsep) * \real{0.3067}}@{}}
\toprule\noalign{}
\begin{minipage}[b]{\linewidth}\raggedright
Переменная
\end{minipage} & \begin{minipage}[b]{\linewidth}\raggedleft
Оценка
\end{minipage} & \begin{minipage}[b]{\linewidth}\raggedleft
Стандартная ошибка (HC3)
\end{minipage} & \begin{minipage}[b]{\linewidth}\raggedleft
Тестовая статистика
\end{minipage} \\
\midrule\noalign{}
\endhead
\bottomrule\noalign{}
\endlastfoot
south & 102.3 & 42.59 & 0.05 \\
\end{longtable}
}

\begin{tcolorbox}[enhanced jigsaw, arc=.35mm, bottomrule=.15mm, bottomtitle=1mm, breakable, colback=white, colbacktitle=quarto-callout-tip-color!10!white, colframe=quarto-callout-tip-color-frame, coltitle=black, left=2mm, leftrule=.75mm, opacityback=0, opacitybacktitle=0.6, rightrule=.15mm, title=\textcolor{quarto-callout-tip-color}{\faLightbulb}\hspace{0.5em}{Вывод}, titlerule=0mm, toprule=.15mm, toptitle=1mm]

\(\lvert t\rvert = 0.05 < 2.58\), гипотеза \(H_0\) не отвергается.

\end{tcolorbox}




\end{document}
